\section{Tool 规格的三层架构}

Codex 源码（主要在 \texttt{spec.rs} 文件中）对工具定义了三层架构：

\subsection{第一层：内部规格（Tool Spec）}

Codex 自己维护的抽象定义，包含：
\begin{itemize}
    \item \textbf{name}：工具名称。
    \item \textbf{description}：工具描述。
    \item \textbf{parameters}：入参定义。
    \item \textbf{output\_schema}：输出结构描述。
\end{itemize}

\begin{warningbox}{output\_schema 不会暴露给 API}
内部规格中定义了 output\_schema（描述工具的返回值结构），但这个字段\textbf{不会}出现在发给远端 API 的请求体中。模型对输出结构的感知，完全依赖于 runtime 将 function call output 回填到上下文中。
\end{warningbox}

\subsection{第二层：Wire Layer（网络传输层）}

"Wire"即变现/网线，指工具在网络传输中的定义。Codex 通过 \texttt{create\_tools\_json\_for\_response\_api} 函数将内部 Tool Spec 序列化为 JSON，塞入 \texttt{response.create} 请求体的 \texttt{tools} 字段。

\begin{knowledgebox}{Function Calling 的设计特点}
标准的 Function Calling 协议中，Tool 定义\textbf{只有入参描述，没有出参描述}。这意味着模型在调用前并不"知道"返回值的精确格式——它通过 runtime 回填的 function call output 来学习返回结构。
\end{knowledgebox}

\subsection{第三层：执行层（Handlers）}

当模型产生一个 Function Call 后，Codex Runtime 进入 Handlers：
\begin{enumerate}
    \item 解析参数并做校验。
    \item 在\textbf{本地}执行操作（执行不在远端服务器上发生）。
    \item 将结果封装为 Function Call Output，放入上下文的消息列表中。
\end{enumerate}

\subsection{本章小结}

三层架构实现了关注点分离：内部规格定义完整的工具语义，Wire Layer 负责序列化传输，Handlers 负责本地执行与结果回填。


